\scanpage{574}

\mspage{483}

Ainsi, la récurrence transfinie se poursuit sur $I$ \deleted{\unsure{tout}} entier. \deleted{Par construction} On voit par récurrence transfinie qu'on a, pour \unsure{tt} $\alpha\in I$

\[s_\alpha\leq x_\alpha\]

(car $x_\alpha>x_\beta\geq s_\beta$ pour \deleted{\unsure{tout}} $\beta<\alpha$, donc $x_\alpha>s_\beta$ pour \unsure{tt} $\beta<\alpha$, \deleted{donc $x_\alpha\geq s_\alpha$} ce qui équivaut à $x_\alpha\geq s_\alpha$). De plus, on trouve \deleted{\unsure{qu'}}\unsure{aussi}, par récurrence \deleted{\unsure{ind}}transfinie, que

\[
\begin{aligned}
\beta<\alpha\ \Longrightarrow\ s_\beta\ &\text{appartient à la sous-catégorie}\\
 &\text{fermée }\inserted{\underline S_\alpha}\text{ de }M\text{ engendrée}\\
 &\text{par les }x_{\beta'}\text{ avec }\beta'<\alpha .
\end{aligned}\tag*{\makebox[0pt][r]{$(*)$\hspace{.90\linewidth}}}
\]

puisque par construction $s_\alpha$ est le \uline{premier} des éléments de $\underline S$ \deleted{tel que} qui n'est \uline{pas} dans \unsure{la dite|ladite} sous-catégorie \deleted{\unsure{de}} $\underline S_\alpha$ de $M$.

Cela montre que \deleted{\unsure{tout}} $s_\beta$ appartient à la sous-catégorie \inserted{fermée} $ {}_{\underline x}M$ de $M$ engendrée par \deleted{\unsure{tous}} \inserted{la famille} de \uline{tous} les \deleted{$\unsure{x}$} $x_\alpha$, \deleted{donc} on trouve que $\underline S\subset{}_{\underline x}M\subset\underline S$ donc $\underline S={}_{\underline x}M$, qed.

\uline{Proposition} 4.16.7.5. Pour un univers $\unsure{\mathcal U|U}$ donné (contenant $\unsure{\mathbb N|N}$ \inserted{comme élément}), les conditions suivantes sont équivalentes:

1) Il existe une catégorie $\mathcal U$-accessible $M$, \deleted{\unsure{et}}\inserted{\unsure{ayant}} une \deleted{sous-}catégorie fermée non accessible

% archive sheet number: 573
